% ============================================================================
%  P1 -- "The structural hypotheses that bound data assimilation"
%  Target venue: JAMES (AGU).
%
%  This file compiles standalone with `article`. To switch to the AGU class on
%  Overleaf, replace the \documentclass line with
%      \documentclass{agujournal2019}
%  and move \title/\authors into AGU's own macros. Everything else is portable.
%
%  SOURCE OF TRUTH: docs/scoping/da_paper_structural_hypotheses.md (SCOPING v2).
%  Every number below is traceable to this repo's own runs; the provenance of
%  each table is given in its caption. Do not edit numbers here without
%  re-deriving them from reports/l96/outputs/ or reports/qg/outputs/.
% ============================================================================
\documentclass[11pt,a4paper]{article}

\usepackage[utf8]{inputenc}
\usepackage[T1]{fontenc}
\usepackage{amsmath,amsfonts,amssymb}
\usepackage{booktabs}
\usepackage{multirow}
\usepackage[margin=2.5cm]{geometry}
\usepackage[dvipsnames]{xcolor}
\usepackage{hyperref}
\usepackage{enumitem}
\usepackage{graphicx}

\hypersetup{colorlinks=true,linkcolor=NavyBlue,citecolor=ForestGreen,urlcolor=NavyBlue}

% --- drafting macros ---------------------------------------------------------
% \todo{...}     work still to do before submission
% \needsrun{...} a claim that is BLOCKED on an experiment not yet run
% \caveat{...}   a claim that is evidenced but carries a stated limitation
\newcommand{\todo}[1]{\textcolor{BrickRed}{\textbf{[TODO: #1]}}}
\newcommand{\needsrun}[1]{\textcolor{BrickRed}{\textbf{[NEEDS RUN: #1]}}}
\newcommand{\caveat}[1]{\textcolor{Sepia}{\textbf{[CAVEAT: #1]}}}

% --- notation ----------------------------------------------------------------
\newcommand{\xv}{\mathbf{x}}
\newcommand{\zv}{\mathbf{z}}
\newcommand{\thv}{\boldsymbol{\theta}}
\newcommand{\Cv}{\mathbf{C}}
\newcommand{\epsv}{\boldsymbol{\epsilon}}
\newcommand{\yv}{\mathbf{y}}
\newcommand{\Hop}{\mathcal{H}}
\newcommand{\Mop}{\mathcal{M}}
\newcommand{\Psiop}{\widetilde{\Psi}}
\newcommand{\Psimean}{\Psiop_{\mathrm{mean}}}
\newcommand{\PsiG}{\Psiop_{\mathrm{G}}}
\newcommand{\PsiNG}{\Psiop_{\mathrm{NG}}}
\newcommand{\Sz}{\textsf{S0}}
\newcommand{\So}{\textsf{S1}}

\title{The structural hypotheses that bound data assimilation:\\
       why more observations stop helping under model error}
\author{Ronan Fablet \\[0.3em]
        \normalsize \todo{affiliation; co-author list}}
\date{Draft --- \today}

\begin{document}
\maketitle

\begin{abstract}
Data assimilation rests on a small number of commitments that are rarely written
down as hypotheses: that the joint inference problem factorises into a sequence
of simpler ones; that skill is bounded by, and increasing in, the fidelity of the
dynamical model; that the state is a point in the phase space of a differential
equation; and that prior errors are propagated by the same integration that
propagates the signal. Each buys something real. We argue that together they
explain two persistent and widely observed symptoms: that data assimilation
converts additional sparse observations into skill inefficiently, and that its
posterior characterisation is structurally impoverished rather than merely
mis-tuned.
%
Using two idealised but complete testbeds --- a two-scale Lorenz-96 system and a
two-layer quasi-geostrophic channel --- we show that model error does not simply
raise the error floor: it \emph{decouples the analysis from the observations}.
Under a $\pm$10\% parameter bias, doubling the observation count buys
strong-constraint 4D-Var a 3.2\% RMSE reduction against 19.7\% in the
perfect-model case, a $6.2\times$ collapse in the marginal value of data, while
gradient-free ensemble filters lose a factor of only $1.9$.
%
We further show that the sensitivity of the free-running model to parameter
perturbations is \emph{not} a proxy for the sensitivity of the posterior to the
same perturbations; that identical physics expressed in a different state
variable produces large and predictable skill differences; and that the
ensemble-Kalman class carries a non-Gaussian deficit by construction rather than
by tuning.
%
This is a diagnosis paper. We do not propose a new assimilation method; we
measure what the standard ones cannot do however well tuned, and attribute each
deficit to a named hypothesis.
\end{abstract}

\section{Introduction}
\label{sec:intro}

Operational data assimilation (DA) is built on commitments so standard that they
are seldom stated as hypotheses at all. The analysis is produced by cycling:
propagate a state estimate, update it when observations arrive, repeat. Skill is
understood to follow model fidelity, so model development and DA development are
the same programme. The state is a point in the phase space of a differential
equation, and its uncertainty is a perturbation of that point --- a covariance,
or an ensemble. Prior errors travel by the same integration that carries the
signal.

Each of these buys something real, and none is foolish. Together, we argue, they
also explain two symptoms that the field observes repeatedly and treats as
tuning problems:

\begin{description}[leftmargin=2em,style=nextline]
  \item[\textnormal{\textbf{S-a --- the marginal value of observations.}}]
    DA converts additional sparse observations into skill inefficiently, and the
    inefficiency worsens sharply under model error. The quantity of interest is
    not RMSE at a given observation density but
    $\mathrm{d}\,\text{skill}/\mathrm{d}\,\text{density}$, measured at matched
    everything.
  \item[\textnormal{\textbf{S-b --- posterior quality.}}]
    DA's posterior characterisation is structurally impoverished, not merely
    mis-calibrated: spread/skill ratios miss their reliability target in ways
    that no inflation parameter repairs, because the deficit is a property of
    the estimator class.
\end{description}

\paragraph{What this paper is, and is not.}
This is a \emph{diagnosis} paper. We propose no new assimilation method. That is
a deliberate choice rather than a limitation of scope: a method paper invites
the reply ``your baseline was under-tuned'', whereas the claim here concerns
what the standard baselines cannot do \emph{however well tuned}. We therefore
report the tuning-sensitivity studies (Section~\ref{sec:tuning}) as evidence
that the baselines were mapped rather than assumed, including one that returned
a negative result.

\paragraph{Contributions.}
\begin{enumerate}[leftmargin=2em]
  \item We state four structural hypotheses (\textbf{H1}--\textbf{H3b},
    Section~\ref{sec:hypotheses}) in a form that admits falsification, and show
    that each can be relaxed \emph{separately} so that the resulting change in
    skill can be attributed to a named commitment.
  \item We show that model error \emph{collapses the marginal value of
    observations} rather than adding a constant error offset
    (Section~\ref{sec:marginal}): a $6.2\times$ collapse for strong-constraint
    4D-Var against $1.9\times$ for ensemble filters. To the best of our survey
    this effect is unreported.
  \item We show that the ODE representation's sensitivity to model parameters is
    \emph{not} a proxy for the posterior's sensitivity to the same parameters
    (Section~\ref{sec:sensitivity}), and that classical DA conflates the two by
    construction because its estimator \emph{is} the model. The consequence is
    methodological: assessing parameter uncertainty by propagating it through
    the model systematically overstates its effect on the posterior whenever the
    observations are informative.
  \item We show that the choice of state variable, at bit-identical physics,
    changes analysis skill substantially and predictably
    (Section~\ref{sec:representation}).
  \item We record a corollary about gradient-based assimilation
    (Section~\ref{sec:adjoint}): the value of the adjoint is conditional on
    model fidelity. At matched conditions the adjoint method and the best
    ensemble filter convert extra observations into skill at an identical rate;
    under model error the adjoint method's return collapses $3.3\times$ harder.
\end{enumerate}

\paragraph{Testbeds.}
All figures are from this study's own runs on two systems: a two-scale
Lorenz-96 model and a two-layer quasi-geostrophic (QG) channel. Both are
idealised. We take that as the price of being able to define ``matched
everything'' precisely enough for a marginal-value measurement to mean
something, and we return to the question of transfer in
Section~\ref{sec:limitations}.

\section{Background and hypotheses}
\label{sec:background}

\subsection{State-space formulation}
\label{sec:statespace}

Let $\xv \in \mathbb{R}^d$ denote the spatiotemporal physical state and
$\zv \in \mathbb{R}^m$ an external forcing. Following the flow-matching
convention we write $\xv_1$ for the \emph{target} state, reserving $\xv_0$ for a
base sample and $\xv_\tau$ for the interpolant introduced in
Section~\ref{sec:methods}. The discrete-time dynamics are
\begin{equation}
  \label{eq:dynamics}
  \xv_{1,k+1} \mid \zv \;=\; \Mop_{\Delta}\!\left(\xv_{1,k}, \zv_{k}, \thv\right) + \xi_k ,
  \qquad k \in \{0,\dots,K\},
\end{equation}
with $\Mop_{\Delta}$ the time-stepping operator over a physical step $\Delta$,
$\thv$ the physical parameters, and $\xi_k$ a state-dependent noise process.
Observations are gathered from the state alone,
\begin{equation}
  \label{eq:obs}
  \yv_k = \Hop(\xv_{1,k}) + \epsv_k , \qquad \epsv_k \sim \mathcal{N}(0, \mathbf{R}).
\end{equation}

Writing $\Phi(\xv_1, \zv; \thv)(t_k) = \Mop_\Delta(\xv_1(t_{k-1}), \zv(t_{k-1}); \thv)$
for $k>0$ and $\Phi(\cdot)(t_0) = \xv_1(t_0)$, the Bayesian state-space model is
\begin{equation}
  \label{eq:ssm}
  \left\{
  \begin{array}{rcl}
    \zv &\sim& \pi_{\zv}, \qquad \thv \;\sim\; \pi_{\thv} \\[0.35em]
    \xv_1 &\sim& \pi_{\xv_1}, \quad
      \pi_{\xv_1}\!\left(\xv_1 \mid \zv, \thv\right) \propto
      \exp\!\left[-\tfrac{1}{2}\left\|\xv_1 - \Phi(\xv_1, \zv; \thv)\right\|^2_{\mathbf{Q}}\right] \\[0.35em]
    \yv &=& \Hop(\xv_1) + \epsv_{\mathrm{obs}}, \quad \epsv_{\mathrm{obs}} \sim \mathcal{N}(0,\mathbf{R}),
  \end{array}\right.
\end{equation}
where $\mathbf{Q}$ is the model-error covariance. The prior on the state is thus
\emph{generated by the model}: $\pi_{\xv_1}$ is defined through $\Phi$, and there
is no other route by which prior information enters. This is the formal content
of H2 and H3b below.

\paragraph{Context, and what this paper estimates.}
In operational practice neither the forcing nor the parameters are known
exactly; the system has access to corrupted surrogates $\zv^{*}$ and $\thv^{*}$.
We collect everything the estimator may condition on into the \emph{context}
\begin{equation}
  \label{eq:context}
  \Cv = \left(\yv_{1:K},\; \zv^{*},\; \thv^{*}\right),
  \qquad
  \Cv_{\leq k} = \left(\yv_{1:k},\; \zv^{*},\; \thv^{*}\right).
\end{equation}
\textbf{Throughout this paper the inference target is the state alone}: we study
the posterior $p(\xv_1 \mid \Cv)$ and its filtered counterpart
$p(\xv_1 \mid \Cv_{\leq k})$. Joint state--parameter inference --- sampling
$p(\xv_1, \thv \mid \Cv)$ --- is a different problem with a different cost, and
is out of scope. Where $\thv^{*}$ appears it is an \emph{input} the estimator may
use, never a quantity it estimates.

\subsection{Four structural hypotheses}
\label{sec:hypotheses}

Classical DA rests on commitments rarely written down as hypotheses.

\begin{description}[leftmargin=1.5em,style=nextline]
  \item[\textbf{H1 (Markovianity).}]
    The joint inference problem over a window --- or a whole reanalysis --- is
    replaced by a \emph{sequence of simpler problems}, each conditionally
    independent of the past given the current state. When $\zv^{*}=\zv$ and
    $\thv^{*}=\thv$, the posterior factorises,
    \begin{equation}
      \label{eq:markov}
      p\!\left(\xv_{1,k'\geq k} \mid \yv_{1:k}, \zv, \thv\right)
      = p\!\left(\xv_{1,k'>k} \mid \xv_{1,k}, \zv, \thv\right)\cdot
        p\!\left(\xv_{1,k} \mid \yv_{1:k}, \zv, \thv\right),
    \end{equation}
    and cycling is exact. This is what licenses filtering; 4D-Var and smoothers
    widen the unit to one window but keep the recursion between windows.
    \emph{The factorisation~\eqref{eq:markov} holds only under those idealised
    conditions.} With $\zv^{*}\neq\zv$ or $\thv^{*}\neq\thv$ the conditional
    independence collapses and the memory induced by the corrupted forcing makes
    the process non-Markovian.
  \item[\textbf{H2 (model fidelity is the route to skill).}]
    The better the dynamical model, the better the assimilation.
  \item[\textbf{H3a (ODE/PDE state representation).}]
    The state is a point in the phase space of~\eqref{eq:dynamics}; uncertainty
    is a perturbation of that point --- a covariance, or an ensemble.
  \item[\textbf{H3b (autoregressive error propagation).}]
    Because $\pi_{\xv_1}$ is generated by integrating $\Phi$, every prior error
    is carried forward by the same recursion that carries the signal, and
    compounds at the system's Lyapunov rate.
\end{description}

\begin{table}[t]
\centering
\small
\begin{tabular}{@{}p{0.06\textwidth}p{0.26\textwidth}p{0.24\textwidth}p{0.32\textwidth}@{}}
\toprule
& \textbf{asserts} & \textbf{buys} & \textbf{costs} \\
\midrule
\textbf{H1} & the joint problem factorises as in~\eqref{eq:markov} & $O(1)$ memory,
operational streaming, no training set; robustness to a change of observing
system without retraining & cannot learn from the \emph{observation
configuration}; no amortisation across windows; the factorisation fails exactly
when $\zv^{*}\neq\zv$ \\[0.4em]
\textbf{H2} & skill is bounded by, and increasing in, model fidelity & a
physically consistent, transferable prior & the model \emph{is} the prior, so
model error enters with no buffer \\[0.4em]
\textbf{H3a} & state $=$ point in phase space; uncertainty $=$ perturbation &
tractable propagation; interpretability & the choice of state variable fixes the
conditioning of the inverse problem; a Gaussian update implies
$\PsiNG \equiv 0$ \\[0.4em]
\textbf{H3b} & prior errors propagate by the same integration as the signal &
one mechanism for signal and uncertainty alike & error compounds at the Lyapunov
rate, and \emph{model} sensitivity is inherited by the analysis whether or not
the \emph{posterior} has it \\
\bottomrule
\end{tabular}
\caption{The four hypotheses, what each buys, and what each costs.}
\label{tab:hypotheses}
\end{table}

\paragraph{H2 is load-bearing, and not simply false.}
Better models do assimilate better. The claim is that H2 is \emph{binding}:
because the model supplies the entire prior via~\eqref{eq:ssm}, DA has no
mechanism to absorb model discrepancy. Model-error covariances, inflation and
weak-constraint formulations exist precisely because this is understood; the
claim is that the available buffers are weak, which
Sections~\ref{sec:inert}--\ref{sec:marginal} measure rather than assert.

\paragraph{H3b is what makes H2 bind.}
The two are not independent. Section~\ref{sec:sensitivity} shows that the
apparent force of ``better model $\Rightarrow$ better DA'' is partly an artefact
of the representation: it is the autoregressive propagation that routes all prior
information through $\Phi$, and so makes the analysis inherit the model's
Lyapunov-amplified sensitivity whether or not the posterior has it.

\section{Conditional flow matching as a generic representation of posteriors}
\label{sec:methods}

This section introduces the instrument used throughout the paper. We do
\emph{not} propose a new assimilation scheme. We use the conditional-flow
representation because it is generic enough to contain every method compared
here as a special case, which turns ``method A is better than method B'' into
the sharper question of \emph{which components of the posterior each method is
able to represent at all}.

\subsection{The conditional flow}

Let $\tau \in [0,1]$ be an abstract transport variable, distinct from physical
time $t$. Define a linear probability path from a Gaussian base sample
$\xv_0 \sim \mathcal{N}(0,\sigma_0^2 \mathbb{I})$ to a target
$\xv_1 \sim p(\xv_1 \mid \Cv)$:
\begin{equation}
  \label{eq:interpolant}
  \xv_\tau = \alpha_\tau \xv_0 + \beta_\tau \xv_1 ,
  \qquad \alpha_0 = 1,\ \beta_0 = 0,\ \alpha_1 = 0,\ \beta_1 = 1 ,
\end{equation}
with the optimal-transport schedule $\alpha_\tau = 1-\tau$,
$\beta_\tau = \tau$. Differentiating~\eqref{eq:interpolant} gives an ODE whose
flow generates the target posterior:
\begin{equation}
  \label{eq:ode}
  \frac{\mathrm{d}\xv_\tau}{\mathrm{d}\tau}
  = \frac{\dot\alpha_\tau}{\alpha_\tau}\,\xv_\tau
  + \left(\dot\beta_\tau - \beta_\tau\frac{\dot\alpha_\tau}{\alpha_\tau}\right)
    \Psiop\!\left(\xv_\tau, \Cv, \tau\right),
  \qquad
  \Psiop\!\left(\xv_\tau, \Cv, \tau\right) \;\triangleq\; \mathbb{E}\!\left[\xv_1 \mid \xv_\tau, \Cv\right].
\end{equation}

Everything method-specific is carried by the single operator $\Psiop$: the
posterior-mean estimate of the target given a partially noised state and the
context. \textbf{$\Psiop$ is the object this paper measures.}

\subsection{The partition}
\label{sec:partition}

$\Psiop$ admits an exact additive decomposition,
\begin{equation}
  \label{eq:partition}
  \Psiop\!\left(\xv_\tau, \Cv, \tau\right)
  = \underbrace{\Psimean(\Cv)}_{\text{conditional mean}}
  + \underbrace{\PsiG(\xv_\tau, \Cv, \tau)}_{\text{affine gain}}
  + \underbrace{\PsiNG(\xv_\tau, \Cv, \tau)}_{\text{non-Gaussian residual}},
\end{equation}
with $\Psimean(\Cv) = \mathbb{E}[\xv_1 \mid \Cv]$ independent of $\xv_\tau$,
$\PsiG$ the $L^2$ projection of $\Psiop - \Psimean$ onto functions affine in
$\xv_\tau$, and $\PsiNG$ the remainder. Defining $\PsiG$ as an $L^2$ projection
makes the three terms orthogonal, so the decomposition is unique and the
attribution in Section~\ref{sec:results} is well posed.

\paragraph{The non-Gaussian branch vanishes at both ends.}
At $\tau=0$, $\xv_\tau \perp \xv_1$, so $\Psiop = \mathbb{E}[\xv_1\mid\Cv] = \Psimean$
and the affine gain is zero; hence $\PsiNG(0)=0$. At $\tau=1$, $\xv_\tau = \xv_1$
and $\Psimean + \PsiG = \xv_1$ exactly; hence $\PsiNG(1)=0$. The non-Gaussian
term is therefore a \emph{bump} on $(0,1)$, pinned to zero at both boundaries.

\subsection{What each method class can represent}
\label{sec:zeroing}

The value of~\eqref{eq:partition} as a diagnostic is that the classical schemes
are not merely different points in a hypothesis space --- they \emph{identically
zero} named terms.

\begin{table}[t]
\centering
\small
\begin{tabular}{@{}lccc@{}}
\toprule
\textbf{method class} & $\Psimean$ & $\PsiG$ & $\PsiNG$ \\
\midrule
Optimal interpolation, 4D-Var & $\hat{\xv}(\Cv)$, constant in $\xv_\tau$ & $\mathbf{0}$ & $\mathbf{0}$ \\
DirectUNet (deterministic regression) & $\hat{\xv}(\Cv)$, constant in $\xv_\tau$ & $\mathbf{0}$ & $\mathbf{0}$ \\
EnKF / ETKF / ensemble smoothers & ensemble mean & \checkmark\ (the Kalman gain \emph{is} $K_\tau$) & $\mathbf{0}$ \emph{by construction} \\
CFM, score-based assimilation (SDA) & \checkmark & \checkmark & \checkmark \\
DirectUNet\,$+$\,SDA (blending) & imported, deterministic & \checkmark & \checkmark \\
\bottomrule
\end{tabular}
\caption{What each class of scheme can represent. The zeros are structural, not
symptoms of tuning: a Gaussian update \emph{is} the ensemble-Kalman method, and a
deterministic estimator ignores its state argument by definition.}
\label{tab:zeroing}
\end{table}

Two consequences frame the rest of the paper.

\paragraph{A deterministic estimator is a degenerate flow.}
Read any deterministic estimator $\hat{\xv}(\Cv)$ as an operator that ignores its
state argument, $\Psiop_{\det}(\xv_\tau,\Cv) = \hat{\xv}(\Cv)$. Substituting
into~\eqref{eq:ode} and integrating gives exactly
$\xv_\tau = (1-\tau)\,\xv_0 + \tau\,\hat{\xv}$. 4D-Var, optimal interpolation and
a supervised regression network are therefore \emph{the same member} of this
family, differing only in how $\hat{\xv}$ is computed. This is why the paper can
place them on one axis without apology.

\paragraph{The ensemble-Kalman class has a structural ceiling.}
$\PsiNG \equiv 0$ is not a deficiency of a particular implementation or
inflation setting: it is what it means to be an ensemble-Kalman method. Any
posterior feature living in $\PsiNG$ is therefore unreachable for that class at
\emph{any} tuning --- which is the precise form of symptom S-b, and what
Section~\ref{sec:nongaussian} measures.

\paragraph{Blending composes the slots.}
Warm-starting a generative sampler from a deterministic estimate --- the
DirectUNet\,$+$\,SDA schemes evaluated here --- is not an unprincipled trick. In
this representation it is the statement that $\Psimean$ may be supplied by one
scheme while $\PsiG$ and $\PsiNG$ are supplied by another. Section~\ref{sec:results}
measures what that composition buys and what it costs.

\section{Evaluation framework}
\label{sec:evaluation}

\subsection{Three case studies, three jobs}
\label{sec:testbeds}

The three systems are not a repeated benchmark; each answers a question the
others cannot.

\begin{description}[leftmargin=2em,style=nextline]
  \item[\textnormal{\textbf{Lorenz-63 --- the model-error axis, including the
    formulation that relaxes H2.}}]
    A stochastic Lorenz-63 system augmented with a wind-forcing variable. Small
    enough that weak-constraint 4D-Var, which explicitly frees a per-step model
    error term, is affordable to run. This matters because weak-constraint
    4D-Var is the classical method \emph{allowed} to know about model error, and
    is therefore the honest ceiling against which H2's cost must be measured.
  \item[\textnormal{\textbf{Lorenz-96 --- the main testbed.}}]
    Two-scale Lorenz-96: eight slow variables coupled to fast variables,
    integrated with RK4 over 500-step assimilation windows. Observations cover a
    24-dimensional subspace of a 40-dimensional state. Every window draws its own
    physical parameters, randomised by $\pm 20\%$, so no method can memorise a
    single setting. This is the only system carrying the full scheme matrix, and
    all cross-class comparisons are made here.
  \item[\textnormal{\textbf{Quasi-geostrophic channel --- representation and
    realism.}}]
    A two-layer Phillips-type channel integrated in potential vorticity $q$ with
    a spectral inversion for the streamfunction $\psi$, forced by a moving-storm
    wind-stress curl. Observations are sparse meridional columns. Its role is
    H3a: the same physics admits two state variables, which is the cleanest
    available test that the representation is not a neutral container.
\end{description}

\begin{table}[t]
\centering
\small
\begin{tabular}{@{}lccccc@{}}
\toprule
& \textbf{DA baselines} & \textbf{DirectUNet} & \textbf{CFM} & \textbf{SDA} & \textbf{DirectUNet$+$SDA} \\
\midrule
Lorenz-63 & \checkmark\ (incl.\ weak-4D-Var) & --- & --- & --- & --- \\
Lorenz-96 & \checkmark & \checkmark & \checkmark & \checkmark & \checkmark \\
QG        & \checkmark\ (EnKF/ETKF) & \checkmark & --- & --- & --- \\
\bottomrule
\end{tabular}
\caption{Scheme coverage. The matrix is deliberately uneven: cross-class claims
are made on Lorenz-96, where every class is present, and the other two systems
carry the specific arguments described above. \todo{Fill the empty cells, or
state in the final version that they are out of scope rather than pending.}}
\label{tab:coverage}
\end{table}

\subsection{Schemes}
\label{sec:schemes}

Five classes, mapped onto Table~\ref{tab:zeroing}:
\textbf{(i)} classical DA --- strong- and weak-constraint 4D-Var, EnKF, ETKF;
\textbf{(ii)} \textbf{DirectUNet}, a single-pass supervised regression
$\Cv \mapsto \hat{\xv}$, i.e.\ a pure $\Psimean$ estimator;
\textbf{(iii)} \textbf{CFM}, a conditional flow trained on the interpolant
\eqref{eq:interpolant};
\textbf{(iv)} \textbf{SDA}, score-based assimilation with an unconditional or
conditioned prior and observation guidance;
\textbf{(v)} \textbf{DirectUNet\,$+$\,SDA}, blending, where the deterministic
estimate supplies $\Psimean$ and the sampler supplies the remaining slots.

Unrolled variational solvers are deliberately excluded. They are a fourth
operator family whose analysis requires the gradient machinery this paper does
not introduce, and they are treated separately.

\subsection{The model-error axis}
\label{sec:s0s1}

\begin{description}[leftmargin=2em,style=nextline]
  \item[\Sz{} (perfect model).] The assimilating model's parameters equal those
    generating the truth for that window.
  \item[\So{} (model error).] The assimilating model's parameters carry an
    additional $\pm 10\%$ bias. The truth is drawn from the same distribution in
    both scenarios; only $\thv^{*}$ changes.
\end{description}

\paragraph{A caveat that must be stated, not cashed in.}
The \Sz{}/\So{} axis prices model error \emph{only for methods that run a forward
model at inference}. A learned estimator receiving no parameters and integrating
nothing is definitionally unaffected, and its $\So/\Sz \approx 1.00$ is not a
robustness result. Three classes must therefore be kept apart
(Table~\ref{tab:classes}), and a model-free ratio is never tabled as cross-class
robustness evidence.

\begin{table}[t]
\centering
\small
\begin{tabular}{@{}p{0.24\textwidth}p{0.36\textwidth}p{0.32\textwidth}@{}}
\toprule
\textbf{class} & \textbf{what \So{} changes for it} & \textbf{a ratio $\approx 1.00$ means} \\
\midrule
classical DA & integrates a biased forward model; the bias compounds through the
integration & n/a --- these degrade by $1.68$--$1.80\times$ \\
model-free learned (DirectUNet, CFM, unconditional SDA) & nothing: no parameters,
no forward model & \emph{definitional}; carries no information \\
conditioned learned (SDA2/SDA3; QG Q2--Q4) & the conditioning inputs
$(\zv^{*},\thv^{*})$ are biased or corrupted & a genuine robustness result \\
\bottomrule
\end{tabular}
\caption{What \So{} means depends on whether, and how, a scheme consumes model
information. This governs every cross-class statement in the paper.}
\label{tab:classes}
\end{table}

\subsection{Metrics}
RMSE is pooled over all windows and timesteps,
$\sqrt{\overline{(\cdot)^2}}$, with the same convention for every scheme.
Explained variance is reported per field. For ensemble schemes the spread/skill
ratio is compared against the finite-ensemble reliability target
$\sqrt{(N-1)/(N+1)}$, and proper ensemble CRPS / energy scores are computed at a
single uniform convention, $N=30$.

\needsrun{Rank histograms are the primary posterior-diagnostic figure for a DA
readership and are not implemented anywhere in the codebase. This remains the
largest new-code item.}

\section{Results}
\label{sec:results}

\subsection{The price of H2}
\label{sec:price}

\begin{table}[t]
\centering
\small
\begin{tabular}{@{}lrrrc@{}}
\toprule
\textbf{scheme} & \textbf{\Sz{}} & \textbf{\So{}} & $\So/\Sz$ & \textbf{slots occupied} \\
\midrule
Strong-4D-Var & 0.8116 & 1.4617 & \textbf{1.801} & $\Psimean$ \\
ETKF          & 0.8883 & 1.4998 & \textbf{1.688} & $\Psimean + \PsiG$ \\
EnKF          & 0.9131 & 1.5381 & \textbf{1.684} & $\Psimean + \PsiG$ \\
\addlinespace
SDA3          & 0.5365 & 0.5355 & 0.998 & all three \\
DirectUNet-M  & 0.5064 & 0.5072 & 1.002 & $\Psimean$ \\
CFM-M         & 0.4810 & 0.4784 & 0.995 & all three \\
DirectUNet\,$+$\,SDA3 & \textbf{0.4204} & 0.4183 & 0.995 & composed \\
\bottomrule
\end{tabular}
\caption{Lorenz-96, 200 windows, pooled RMSE. \textbf{The two blocks are
different method classes} and the lower ratios must be read through
Table~\ref{tab:classes}: DirectUNet and CFM are model-free, so their
$\So/\Sz \approx 1$ is definitional. Only the classical block's degradation
measures the cost of model error. Source:
\texttt{reports/l96/outputs/l96\_consolidated\_benchmark.md}.}
\label{tab:price}
\end{table}

Classical DA degrades by $1.68$--$1.80\times$ from \Sz{} to \So{}. That is the
unbuffered cost of H2: by~\eqref{eq:ssm} the model \emph{is} the prior, so a
parameter bias is not a perturbed input but a corrupted prior, compounding over
a 500-step window.

\subsection{Weak-constraint 4D-Var, the formulation that relaxes H2}
\label{sec:weak4dvar}

Lorenz-63 is small enough to afford the weak-constraint formulation, which frees
a per-step model-error term and is therefore the classical method permitted to
know that the model is wrong.

\begin{table}[t]
\centering
\small
\begin{tabular}{@{}lrr@{}}
\toprule
\textbf{scheme} & \textbf{\Sz{}} & \textbf{\So{}} \\
\midrule
Weak-4D-Var   & \textbf{0.642} & \textbf{1.633} \\
Strong-4D-Var & 0.767 & 2.100 \\
EnKF          & 0.881 & 2.267 \\
ETKF          & 0.878 & 2.271 \\
\bottomrule
\end{tabular}
\caption{Lorenz-63 state RMSE. Weak-4D-Var is best in both scenarios, and under
model error it beats strong-constraint by 22\%.
Source: \texttt{reports/l63/outputs/s0\_s1\_synthesis.md}.}
\label{tab:weak4dvar}
\end{table}

This is the honest ceiling for the classical family, and it sharpens rather than
weakens the argument. Relaxing H2 \emph{within} the classical framework buys a
great deal --- so the residual gap to the learned schemes is not an artefact of
comparing against a strawman. It also bounds what C1 may claim: the
marginal-value collapse reported in Section~\ref{sec:marginal} is established for
strong-constraint and filtering methods, and the corresponding weak-constraint
measurement on Lorenz-96 does not yet exist.
\needsrun{Weak-constraint 4D-Var on Lorenz-96 across the observation-density
grid. The implementation exists; the benchmark run does not.}

\subsection{Conditioning on model information is nearly inert}
\label{sec:inert}

If the learned schemes are robust to \So{}, the obvious question is whether they
use $\thv^{*}$ at all.

\begin{table}[t]
\centering
\small
\begin{tabular}{@{}llrrr@{}}
\toprule
\textbf{scheme} & \textbf{conditioning} & \textbf{\Sz{}} & \textbf{\So{}} & $\So/\Sz$ \\
\midrule
SDA1 & none (unconditional prior) & 0.5532 & 0.5521 & 0.998 \\
SDA2 & \emph{true} $\thv$ $+$ corrupted $\zv^{*}$ & 0.5588 & 0.5610 & 1.004 \\
SDA3 & \emph{noisy} $\thv^{*}$ & \textbf{0.5365} & 0.5355 & 0.998 \\
\bottomrule
\end{tabular}
\caption{Lorenz-96 conditioning ladder, pooled RMSE. Conditioning on \emph{true}
parameters is 1.0\% \emph{worse} than not conditioning at all; conditioning on
\emph{noisy} ones is 3.0\% better.}
\label{tab:ladder}
\end{table}

Across three independent families the effect of model-information conditioning
on \Sz{} skill is at most $\sim$1.5\%, and \emph{the sign is not consistent}: the
SDA family gives $+1.0\%$ (worse) for true parameters and $-3.0\%$ (better) for
noisy ones, a non-monai SDA variant gives $-1.5\%$, and a $\tau{=}0$ CFM
conditioned on corrupted forcing gives $+0.9\%$ (worse). Conditioning is, to a
good approximation, \emph{inert}; the one variant that reliably helps is the
noisy one, which looks like regularisation rather than information use.

This is the sharpest available statement of the thesis. The learned schemes are
robust to corrupted model parameters largely \emph{because they barely depend on
them} --- the observations carry the information. Classical DA has no such
fallback: its prior \emph{is} the integrated model. H2 binds not because model
information is unhelpful, but because DA cannot down-weight it.

\paragraph{The training-exposure control has been run.}
The configuration trained with zero forcing/state bias --- the only one never
exposed to \So{}-level error in training --- gives \Sz{} 0.7046 / \So{} 0.7033,
ratio 0.998, indistinguishable from its exposed sibling (0.997). Removing
training-time exposure entirely leaves the invariance intact.
\caveat{This control has no counterpart in the canonical benchmark table, which
is scoped to one backbone family; it is cited to its own run artefact.}

\paragraph{The same ladder on QG.} The QG DirectUNet variants repeat the pattern
at a different scale and resolution: the oracle-conditioned variant attains the
best \Sz{} explained variance (0.9827) but falls to 0.9490 under \So{}, while the
noisy-conditioned variants hold (0.9798 $\to$ 0.9763, and 0.9824 $\to$ 0.9808).
Conditioning on \emph{true} model information is again the fragile choice.

\subsection{ODE sensitivity is not posterior sensitivity}
\label{sec:sensitivity}

Two distinct sensitivities are in play, and classical DA cannot separate them:
the \emph{model} sensitivity $\partial \xv_{\mathrm{ODE}}/\partial(\thv,\zv)$ of
the free-running trajectory, Lyapunov-amplified and large; and the
\emph{posterior} sensitivity
$\partial\, p(\xv_1\mid\Cv)/\partial(\thv,\zv)$ at fixed $\yv$.

Classical DA conflates them \textbf{by construction}, because its estimator
\emph{is} the model: by~\eqref{eq:ssm} every route from $(\thv,\zv)$ to the
analysis passes through $\Phi$. Its \Sz{}$\to$\So{} degradation therefore measures
\emph{its own estimator's} sensitivity, and cannot be read as a measurement of
how much $(\thv,\zv)$ uncertainty matters for the posterior.

A conditioned amortised estimator measures the other quantity directly: it takes
$(\thv^{*},\zv^{*})$ as inputs at fixed $\yv$, so perturbing the conditioning
while holding observations fixed is a finite-difference estimate of posterior
sensitivity. That perturbation moves skill by $\leq 1.5\%$
(Section~\ref{sec:inert}), and full corruption by $\sim 0.2\%$.

\paragraph{Consequence.}
Assessing parameter and forcing uncertainty by \emph{propagating it through the
model} systematically \textbf{overstates} its effect on the posterior whenever
the observations are informative, and the overstatement is exactly the gap
between the two sensitivities --- here $1.68$--$1.80\times$ against $\leq 1.5\%$.

\subsection{Model error destroys the marginal value of observations}
\label{sec:marginal}

\begin{table}[t]
\centering
\small
\begin{tabular}{@{}lccc@{}}
\toprule
\textbf{method} & \textbf{\Sz{}: gain from $2\times$ obs} & \textbf{\So{}: same gain} & \textbf{collapse} \\
\midrule
Strong-4D-Var & $0.9701 \to 0.7788$ (\textbf{19.7\%}) & $1.4751 \to 1.4276$ (\textbf{3.2\%}) & $\mathbf{6.2\times}$ \\
ETKF          & $1.0973 \to 0.8815$ (19.7\%)          & $1.6367 \to 1.4680$ (10.3\%)         & $1.9\times$ \\
EnKF          & $1.0927 \to 0.9046$ (17.2\%)          & $1.6503 \to 1.5022$ (9.0\%)          & $1.9\times$ \\
\bottomrule
\end{tabular}
\caption{Lorenz-96, \emph{temporal} observation density doubled (15 $\to$ 30
observations per window). A separate configuration from Table~\ref{tab:price}
--- all five parameters randomised $\pm 20\%$ --- so the two must not be
cross-quoted. Source:
\texttt{reports/l96/outputs/s0\_s1\_obs\_density\_da\_baselines.md}.}
\label{tab:marginal}
\end{table}

Under model error, doubling the observations buys strong-constraint 4D-Var
3.2\%. In the fast-variable subspace, explained variance goes
$+0.417 \to -0.011$ at the sparser density and $+0.622 \to +0.047$ at the denser
one.

\textbf{H2 does not merely set an error floor; it decouples the analysis from the
observations.} A method whose prior is a wrong model cannot be rescued by more
data --- exactly the regime of sparse-observation geophysical DA. This is also
H3b's shadow: at low density DA has nothing but the model, so model error and
observation sparsity \emph{compound rather than add}.

\subsection{The state variable changes skill at identical physics}
\label{sec:representation}

The QG \texttt{psi\_state} variant integrates the same dynamics with the
streamfunction as the state variable. The $q$-space physics is bit-identical ---
free forecasts agree to $\sim 10^{-6}$ relative --- and only the representation
differs, and with it the observation operator, which becomes a trivial index
lookup.

It gives the \emph{best per-field streamfunction analysis}
($\psi_1,\psi_2$ EV $\approx 0.98$) and a \emph{degenerate potential-vorticity
field} ($q$ EV $=-3.2$ against the $q$-state reference $+0.34$), because
$q \approx \nabla^2\psi$ amplifies high-wavenumber $\psi$-analysis error by
$K^2$. At matched observation noise the $\psi$-state $q$-field EV is 0.583
against the $q$-state 0.752.

Identical dynamics, identical information, and a large skill difference produced
purely by which variable is called ``the state''. The representation is not a
neutral container.
\caveat{The headline $-3.2$ is partly a metric choice --- the aggregate score
reports the $q$ field. The per-field numbers above are the ones to quote.}

\subsection{The non-Gaussian deficit, and what blending buys}
\label{sec:nongaussian}

Table~\ref{tab:zeroing} predicts that the ensemble-Kalman class cannot represent
$\PsiNG$ at any tuning. The measured consequence is that \emph{no scheme
evaluated is simultaneously accurate and calibrated}. The only configuration
carrying a full non-Gaussian branch is 20--25\% over-dispersive --- spread/skill
relative to target $1.20 / 1.22 / 1.24$ at $N = 6/10/30$, stable in $N$ and hence
a sampler property rather than a sample-size artefact --- while every scheme that
closes the RMSE gap does so by importing a deterministic $\Psimean$ and collapses
to ratios $0.28$--$0.59$ against a target of $0.967$.

Blending is where this becomes a design choice rather than an accident.
DirectUNet\,$+$\,SDA3 is the most accurate scheme in Table~\ref{tab:price}
(0.4204, against 0.5064 for its own $\Psimean$ alone and 0.5365 for the sampler
alone), so composing the slots demonstrably buys accuracy. But it buys it by
importing a deterministic mean, which is precisely the move that collapses
dispersion. \textbf{Accuracy and calibration are being traded against each other
through the same mechanism}, and the partition~\eqref{eq:partition} is what makes
that visible: the trade is not between schemes but between which slots are filled
generatively and which are imported.
\needsrun{A calibrated comparison across all classes at one uniform ensemble
convention, with rank histograms. Several cells of the benchmark are still
single-pass point estimates.}

\subsection{A motivating ordering that is not yet evidence}
\label{sec:ordering}

At \Sz{} on Lorenz-96 the schemes order by how much of the window they treat at
once: sequential filters (ETKF 0.8883, EnKF 0.9131) are worse than the window
variational method (Strong-4D-Var 0.8116), which is in turn far worse than the
window-amortised learned schemes (0.4204--0.5365). Read against
Table~\ref{tab:zeroing}, the ordering is not monotone in the number of operator
slots occupied --- EnKF fills $\PsiG$ and still trails a $\Psimean$-only
variational scheme --- which is itself informative: \emph{how much of the problem
is solved jointly} matters here more than which posterior components are
representable.

This is suggestive of H1 but \textbf{confounded}: amortisation, learning and
capacity all vary at once, and the learned schemes see a training set the
classical ones do not.
\needsrun{Filter vs.\ window-smoother vs.\ window-amortised at matched model,
matched observations and matched compute. This is what turns the ordering into
evidence for H1.}

\section{Discussion}
\label{sec:discussion}

\subsection{Claims}
\label{sec:claims}

\begin{description}[leftmargin=2.2em,style=nextline]
  \item[\textbf{C1 (H2, headline).}]
    Because the model \emph{is} the prior, model error enters the analysis
    unbuffered, and its cost is not a constant offset: for strong-constraint
    4D-Var it \emph{collapses the marginal value of observations} ($6.4\times$;
    $7.0\times$ at the benchmark inflation). The ensemble filters keep most of it
    ($1.1\times$; $1.6$--$1.7\times$ at the benchmark inflation).
    \caveat{Established for strong-constraint 4D-Var; the filters' collapse
    reported in earlier drafts ($1.9\times$) came from a mis-specified \Sz{} and
    does not survive the fix (Section~\ref{sec:marginal}). Lorenz-63
    shows weak-constraint 4D-Var recovers much of the loss
    (Section~\ref{sec:weak4dvar}); the corresponding Lorenz-96 measurement on the
    density grid does not yet exist, so C1 is stated for strong-constraint
    4D-Var and not for the classical framework as a whole.}
  \item[\textbf{C2 (H3a).}]
    The state representation fixes the conditioning of the inverse problem.
    Identical physics in a different state variable produces large, predictable
    skill differences via $q\approx\nabla^2\psi$ and its $K^2$ amplification.
  \item[\textbf{C3 (H1).}]
    Markovian factorisation forecloses \emph{learning} the observation
    configuration, and buys robustness to it for free. DA is the scheme least
    affected by a change of observing system ($\times 1.06$--$1.20$); an amortised
    estimator trained on one configuration is not ($\times 2.7$--$3.3$). Trained on
    the \emph{distribution} of configurations, it loses little
    ($\times 1.24$--$1.32$), matches the fixed-grid model on that grid, and beats DA
    from about 20 observations per window at \Sz{} and at every density at \So{}
    (Section~\ref{sec:ordering}).
  \item[\textbf{C4 (H1 $+$ H3a $\Rightarrow$ S-b) --- \emph{not supported as
    originally stated}.}]
    The intended claim was that the ensemble-Kalman class, carrying
    $\PsiNG \equiv 0$ by construction, has a structurally impoverished posterior.
    \textbf{Neither dispersion nor shape measurements support it}
    (Section~\ref{sec:nongaussian}). ETKF and EnKF are the best-calibrated
    schemes measured (spread/RMSE $0.98$, $0.97$ against a target of $0.967$),
    while every scheme with a full non-Gaussian branch is under-dispersed by
    $1.4\times$ to $2.3\times$. Rank histograms --- which could see a shape
    deficit that spread/RMSE cannot --- show the filters' residual shape error,
    once spread and mean bias are corrected, in the same range as
    PredictStateCFM's and below SDA's. \textbf{C4 is therefore not supported as a
    discriminating claim}, and S-b rests on the accuracy/dispersion trade being
    unresolved by \emph{every} class rather than on a deficit specific to one.
  \item[\textbf{C5 (attribution).}]
    Each hypothesis can be relaxed separately and the resulting change
    attributed to a named term of~\eqref{eq:partition}. This is what makes the
    paper a diagnosis rather than a benchmark.
  \item[\textbf{C6 (H2 $\times$ H3b).}]
    The representation's sensitivity to $(\thv,\zv)$ is \emph{not a proxy} for
    the sensitivity of $p(\xv_1\mid\Cv)$ to $(\thv,\zv)$. DA conflates them
    because its estimator is the model; a conditioned amortised estimator
    separates them; and the gap is large --- $2.0$--$2.2\times$ against $0$--$2\%$
    (Section~\ref{sec:sensitivity}).
\end{description}

\subsection{Limitations and threats to validity}
\label{sec:limitations}

\paragraph{The \Sz{}/\So{} axis means different things per class.}
\So{} is a \emph{no-op} for model-free learned schemes; their $\So/\Sz\approx 1.00$
carries no information and is never tabled as robustness
(Table~\ref{tab:classes}). Even for conditioned schemes, being handed biased
parameters is not the same exposure as \emph{integrating} a biased model for 500
steps: one is an input perturbation, the other compounds. Section~\ref{sec:marginal}'s
intra-classical result needs neither caveat.

\paragraph{``Your baselines were under-tuned.''}
The standard reply to any DA-versus-learned comparison. Our mitigation is
evidential: localisation and inflation sensitivity studies found the QG
localisation radius untuned \emph{even at the reference density}, with an optimum
that moves with observation density --- itself an instance of what H1 forecloses
--- and a 4D-Var variance-scaling sweep returned a \emph{negative} result, the
existing defaults already being best. On Lorenz-96 the objection was partly
right, and we disclose it: the DA forward model had not received the
fast-variable coupling weights and \Sz{} ran at the \So{}-tuned inflation; fixing
both lowered ETKF's \Sz{} RMSE by 21\%. Every number here is the fixed one.
Section~\ref{sec:weak4dvar} is the stronger answer: the classical family is
represented by its best member, not its most convenient one.

\paragraph{Uneven scheme coverage.}
Only Lorenz-96 carries all five classes (Table~\ref{tab:coverage}). Cross-class
claims are made there; the other two systems carry specific arguments. This is a
design choice, but it does mean the generality of the cross-class ordering rests
on one system.

\paragraph{Different configurations must not be cross-quoted.}
Tables~\ref{tab:price} and~\ref{tab:marginal} come from different configurations
--- different observation counts, different randomisation breadth. Ratios from
one are not comparable to levels in the other. The same holds across test sets
(regular grid versus random observing system) and across training protocols
(P1 versus benchmark default, 400 versus 1200 epochs): the learned ranking itself
changes between them.

\paragraph{The circularity objection.}
An amortised estimator needs a simulator to generate its training data. If that
simulator is the truth model, the comparison has handed it the truth --- and in
operations there is no truth model. This is the deepest objection to this line of
work and we do not claim to answer it fully. The framing is therefore about
\emph{where prior information enters and what it costs}, not ``learned beats
physical''.
\todo{A self-supervised objective --- a weak-constraint variational cost
evaluated on the estimator's own output, with no ground-truth supervision --- is
implemented and is a direct, if partial, answer. Not yet run; decide whether it
belongs here or in the companion paper.}

\paragraph{Idealised testbeds.}
All three systems are idealised. That is what makes ``matched everything''
definable precisely enough for a marginal-value measurement to mean anything,
but transfer to operational systems is an open question we do not test.

\subsection{Positioning}
\label{sec:positioning}

\todo{Every citation is from memory and must be verified before a bibliography
exists. Nothing in \texttt{refs.bib} has been checked, and no \texttt{$\backslash$cite}
commands are used yet.}

The DA review literature frames the field's standard commitments; the
assimilation-in-the-unstable-subspace line gives the closest existing account of
\emph{why} DA's effective degrees of freedom are limited; weak-constraint 4D-Var
is the canonical relaxation of H2; innovation-based diagnostics supply the
standard consistency checks, whose assumptions are violated when the estimator is
trained to fit the observations. On the learned side, the
DA-with-learned-model-error-correction line and the score-based and flow-based
assimilation line are the nearest neighbours.

\paragraph{The gap.}
The model-error literature asks how to \emph{correct} the model; the
unstable-subspace literature explains the dimensionality of the analysis; the
learned-DA literature proposes hybrids. As far as our survey goes, none of them
\emph{measures the marginal value of observations as a function of model error},
or attributes the deficit to named components of the inference operator.
Section~\ref{sec:marginal} is, on that reading, an unreported effect.

\subsection{Open experiments}
\label{sec:open}

\todo{This list is provisional. It was drawn up against the previous,
hypothesis-ordered outline; it should be re-derived from the structure above now
that the results lead with the cross-system overview, and pruned to what the
paper actually needs rather than everything that could be measured.}

\begin{description}[leftmargin=2.2em,style=nextline]
  \item[\textbf{Weak-constraint 4D-Var on Lorenz-96.}] Bounds C1. Implemented;
    the benchmark run does not exist. Highest value, lowest cost.
  \item[\textbf{H1 isolated.}] Filter vs.\ window-smoother vs.\ window-amortised
    at matched model, observations and compute --- what turns
    Section~\ref{sec:ordering}'s ordering into evidence.
  \item[\textbf{Model-error dose--response.}] Sweep model error continuously
    rather than the binary \Sz{}/\So{}, giving the \emph{same} biased parameters
    to the classical forward model and to the learned conditioning inputs.
  \item[\textbf{The marginal-value surface.}] Observation density $\times$ model
    error for every class on matched axes. A factorial grid of random observing
    systems (6 observation counts $\times$ 4 fast-channel counts, the same layouts
    for every class) now covers the density axes at \Sz{} and \So{}; the missing
    part is a continuous model-error axis.
  \item[\textbf{The two sensitivities on one grid.}] Model sensitivity and
    posterior sensitivity over observation density $\times$ perturbation
    amplitude; the crossover density is the paper's second deliverable number
    after Section~\ref{sec:marginal}'s collapse factor.
\end{description}

\section{Conclusion}
\label{sec:conclusion}

Data assimilation's standard commitments are not errors. They buy streaming
operation, physical interpretability, tractable propagation, and a prior that
transfers. What we have tried to show is that they also \emph{bind}, in a way
that is measurable and attributable to named terms of a single operator.

The headline is that model error does not simply add an error floor. It
decouples the analysis from the observations: under a modest parameter bias, the
return on doubling the observation count falls by a factor of six for
strong-constraint 4D-Var, whose analysis is constrained to the wrong model's
trajectory; the ensemble filters keep most of it. In the sparse-observation regime that characterises geophysical
assimilation, that is the difference between a system that improves with more
data and one that does not.

Two further results sharpen the mechanism. The sensitivity of the free-running
model to parameter perturbations is not the sensitivity of the posterior to the
same perturbations, and classical DA cannot tell them apart because its
estimator \emph{is} the model --- so the routine practice of pricing parameter
uncertainty by propagating it forward overstates its cost whenever the
observations are informative. And the representation is not a neutral container:
identical physics in a different state variable changes analysis skill
substantially, in a direction predictable from the conditioning of the map
between representations.

Reading these through the conditional-flow representation is what makes the
diagnosis attributable rather than anecdotal: the classical schemes
\emph{identically zero} named components of the posterior-mean operator. The
measurements qualify how much that matters. Occupying more components does not by
itself buy accuracy --- a mean-only learned estimator matches the full conditional
flows --- and the zeroed non-Gaussian component shows up neither in dispersion nor
in rank histograms. What the measurements single out is where model information
enters: the most accurate scheme measured, a deterministic estimate blended into a
generative sampler, is also flat under model error, but only when its prior was
trained on noisy rather than true model parameters.

We have deliberately not proposed a replacement. The value of a diagnosis is
that it states what a class of methods cannot do however well tuned, and the
evidence here is arranged so that each deficit attaches to a named hypothesis
rather than to an implementation.

\todo{Revisit once weak-constraint 4D-Var has run on Lorenz-96. Lorenz-63
already shows it recovering much of the model-error loss, so the conclusion may
need to be narrowed to strong-constraint 4D-Var --- narrower, but
better supported, and worth stating plainly rather than hedging.}

\section*{Open science statement}
All figures derive from simulation code and experiment configurations in the
accompanying repository; no external dataset is used. \todo{Add the archive DOI
and the commit hash used for the final numbers.}


\bibliographystyle{plainnat}
\bibliography{refs}

\end{document}
